% Primitive: ctc-gradient — authored from deep_learning/ctc_gradient_manim.py
% (seven scenes), its README concepts table and plan 010's verification
% pass (exact rational arithmetic, two independent routes; anchors A-T).
% This is the guide's destination: it spends every earlier chapter's
% blocks and points forward only to decoding. Conventions inherited from
% the plan's decisions: the 2012 (exclusive) beta, the log-sensitivity
% derivation (no product rule), the loss-gradient sign with the push
% named where it appears, blank dominance conditioned on T. Problem
% answers are computed by answers/ctc_gradient.py (the solve-gate anchor).
\chapter{The CTC gradient: softmax minus occupancy}

\begin{narrative}
This is what the road was for. The alignment chapter left a loss —
the negative log of a sum over every path that spells the
transcript — and every chapter since has been quietly forging a tool
to differentiate it. The worked example returns one last time:
transcript $Y = \mathrm{AB}$, alphabet $\{A, B, \varepsilon\}$,
$T = 4$ frames, the wrapped states $(\varepsilon, A, \varepsilon, B,
\varepsilon)$, the $\anchor{001.paths15}$ paths — and now a concrete
per-frame matrix (rows $A/B/\varepsilon$, one column per frame):
$\anchor{010.K.matrix}$. On it, exact arithmetic gives
$P(\mathrm{AB} \mid X) = \anchor{010.K.P}$ and loss
$-\ln P = \anchor{010.K.NLL}$. By the chapter's end, the gradient of
that loss at every one of the twelve score cells will be a
subtraction you can read off a picture.

\section{The other half of the trellis}

The forward variable $\alpha_t(s)$ answered ``how did we get
here?''. Its mirror $\beta_t(s)$ answers ``how do we finish?'' — the
weight of every legal way to complete the transcript from state $s$
after frame $t$. Same grid, same edges, swept against the arrows.

One bookkeeping decision matters more than it looks: each frame's
emission is a coin exactly one of the two variables may pocket.
This book adopts the convention of Graves's 2012
treatment~\cite{deep_learning-alex-graves-supervised-sequence-labelling-with-rnns-2012}:
$\alpha$'s column pockets frame $t$'s emission, and $\beta$ starts
its product at $t + 1$ — so at the last frame, $\beta_T = 1$ at the
two accepting states and $0$ elsewhere, and no emission is ever
counted twice. At unit weights the backward sweep is pure counting:
its first column reads $(5, 10, 6, 4, 1)$ down the five states, and
the two \emph{initial} states give $5 + 10 = \anchor{001.paths15}$ —
the flagship count, found now from the far end. (On this example the
backward columns mirror the forward ones — a symmetry of this
palindromic transcript, not a theorem.) Formula last, the mirror
recurrence:
\[
  \beta_t(s) \;=\; \sum_{i} \beta_{t+1}(s{+}i)\; y_{t+1}(z'_{s+i}),
\]
the sum over stay, advance, and the skip where the grid allows it.
This is the backward half of forward--backward — the HMM lineage
every CTC implementation
inherits~\cite{deep_learning-lawrence-r-rabiner-the-hmm-tutorial-proc-ieee-1989}.

\section{Paths through a cell}

Multiply the two halves at any cell and something new appears:
$\alpha_t(s)\,\beta_t(s)$ is the total weight of every path that
passes \emph{through} that cell — ways in, times ways out, the
counting grid's product rule reborn carrying probability. At unit
weights, cell $(t{=}2, A)$ has $2$ prefixes and $4$ suffixes:
$2 \times 4 = 8$ of the $15$ paths cross it.

\begin{center}
\begin{tikzpicture}[x=1.55cm, y=0.95cm]
  % alpha*beta path counts through each cell at unit weights — forward
  % counts (plan 001) times backward counts (plan 010 anchor P). The
  % t=2 column is highlighted: every column sums to the same 15.
  \fill[Accent!18] (1.62, -6.25) rectangle (2.38, -0.45);
  \foreach [count=\t from 1] \col in
    {{5,10,0,0,0},{1,8,3,3,0},{0,3,3,8,1},{0,0,0,10,5}} {
    \foreach [count=\s from 1] \v in \col {
      \ifnum\v>0
        \node[circle, draw=Muted, minimum size=6.4mm, inner sep=0pt]
          (n\t\s) at (\t, -\s) {\small $\v$};
      \else
        \node[circle, draw=Muted!45, minimum size=6.4mm, inner sep=0pt]
          (n\t\s) at (\t, -\s) {};
      \fi
    }
    \node[above, color=CoolInk] at (\t, -0.42) {\small $t=\t$};
    \node[color=AccentInk] at (\t, -5.95) {\small $15$};
  }
  \foreach [count=\s from 1] \lab/\col in
    {$\varepsilon$/Muted, $A$/PaletteA, $\varepsilon$/Muted, $B$/PaletteB, $\varepsilon$/Muted}
    \node[left, color=\col] at (0.6, -\s) {\lab};
  \foreach \t/\u in {1/2, 2/3, 3/4} {
    \foreach \s/\r in {1/1,1/2, 2/2,2/3,2/4, 3/3,3/4, 4/4,4/5, 5/5}
      \draw[Muted!60] (n\t\s) -- (n\u\r);
  }
  \node[right, color=AccentInk] at (4.45, -5.95)
    {\small every column: $\anchor{001.paths15}$};
\end{tikzpicture}
\end{center}

Sweep a column across the grid and the series' strongest device
appears — Eisner's constant
column~\cite{deep_learning-jason-eisner-the-forward-backward-spreadsheet-acl-02}:
at unit weights \emph{every} column of products sums to $15$, and on
the real matrix every column sums to the same
$P(\mathrm{AB} \mid X) = \anchor{010.K.P}$ — four columns, one
answer, which also equals $\alpha$'s two-final-state sum and
$\beta$'s two-initial-state sum. The reason is a sentence: every
path crosses every column exactly once. Formula last:
\[
  P(Y \mid X) \;=\; \sum_s \alpha_t(s)\,\beta_t(s)
  \qquad \text{for any } t .
\]
The column is also the bookkeeping self-check: pocket an emission
twice and its column's sum scales visibly out of line.

\section{Where the truth spends its time}

Divide each column by its own sum and it becomes a distribution:
\[
  \gamma_t(s) \;=\; \frac{\alpha_t(s)\,\beta_t(s)}{P(Y \mid X)} ,
\]
the share of the truth's probability passing through each cell —
\emph{where the truth spends frame $t$}. The divide-by-$P$ is the
conditional chapter's renormalized slice, conditioned this time on
the transcript; columns sum to $1$ because each path occupies
exactly one cell per frame — the topology grants it. Folding blank's
three grid rows into one class (a summation drawn exactly once — the
grid has five rows, the bars three) gives the per-class occupancy on
the real matrix:

\begin{center}
\begin{tabular}{c|cccc}
 & $t=1$ & $t=2$ & $t=3$ & $t=4$ \\
\hline
$A$ & $0.9028$ & $0.7086$ & $0.1464$ & $0$ \\
$B$ & $0$ & $0.0330$ & $0.1403$ & $0.9487$ \\
$\varepsilon$ & $0.0972$ & $0.2584$ & $0.7133$ & $0.0513$ \\
\end{tabular}
\end{center}

Columns are distributions; \emph{rows are not}. A's row sums to
$\anchor{010.S.dwellA}$ — an expected \emph{dwell time}: of the four
frames, the truth spends nearly two of them emitting $A$, and the
three dwells $\anchor{010.S.dwellA} + \anchor{010.S.dwellB} +
\anchor{010.S.dwellblank}$ sum to exactly $4 = T$. The balance point
from the random-variables chapter has returned: occupancy is an
expectation over the posterior path distribution. One degenerate
setting shows what the object is made of: under \emph{uniform}
outputs ($y = 1/3$ everywhere), $P = \anchor{010.L.uniformP}$ and
$\gamma$ collapses to pure path counts over $15$ — the alignment
chapter's counting scene reborn as a distribution. Stand a $\gamma$
column beside the derivative chapter's one-hot target and the
chapter's destination is already visible: this is a target
\emph{gone soft}.

\section{The sensitivity of the sum}

Now differentiate the monster — and find there is no monster. First,
a linearity observation: each path uses frame $t$ exactly once, so
any single cell's score $y_t(k)$ appears in each path's product to
power $0$ or $1$ — $P$ is \emph{linear} in that one cell, and its
slope is the cell's coefficient, which the waist picture already
named: $\alpha \beta$ summed over the cell's states.

The load-bearing step is the log-sensitivity reading, and it is the
derivative chapter's closing move at path scale. Nudge one cell
multiplicatively — scale $y_t(k)$ by $(1 + h)$ — and every path
through the cell scales by the same factor while every other path
stands still. So the \emph{relative} change of $P$ is the share of
$P$ passing through the cell:
\[
  \frac{\partial \ln P}{\partial \ln y_t(k)} \;=\; \gamma_t(k) ,
\]
with the score function $d \ln f = f'/f$ doing the bookkeeping, as
it has since the toolkit chapter. In $\mathrm{LSE}$'s gradient each
term owned its own score; here many paths share one cell, so their
shares \emph{add} into occupancy — the sensitivity of a smooth max
is a share, and the sensitivity of a path-sum is a pooled share.
Since the loss is $L = -\ln P$, the formula lands last:
\[
  \frac{\partial L}{\partial \ln y_t(k)} \;=\; -\gamma_t(k) .
\]
Nudge one cell's log-probability, and the loss moves by exactly that
cell's occupancy.

\section{Softmax minus occupancy}

The per-frame $y$ is a softmax, so one owned step remains. The
log-softmax derivative — one line from
$\partial_i \mathrm{LSE} = \mathrm{softmax}_i$ — is
$\partial \ln y_j / \partial u_k = \delta_{jk} - y_k$, where $u$ are
the frame's logits. Chain it against $-\gamma$:
\[
  \frac{\partial L}{\partial u_t(k)}
  \;=\; -\sum_j \gamma_t(j)\,\bigl(\delta_{jk} - y_t(k)\bigr)
  \;=\; -\gamma_t(k) + y_t(k) \sum_j \gamma_t(j) ,
\]
and the checksum $\sum_j \gamma_t(j) = 1$ — the constant column,
paying its second time, now load-bearing in the algebra — collapses
the whole thing:
\[
  \boxed{\;\frac{\partial L}{\partial u_t(k)}
    \;=\; y_t(k) - \gamma_t(k)\;}
\]
(Graves 2012, eq.\ 7.34~\cite{deep_learning-alex-graves-supervised-sequence-labelling-with-rnns-2012}).
The per-frame gradient of the CTC loss is the softmax output minus
the truth's occupancy. On the worked matrix:

\begin{center}
\begin{tabular}{c|cccc}
$y - \gamma$ & $t=1$ & $t=2$ & $t=3$ & $t=4$ \\
\hline
$A$ & $-0.2028$ & $-0.1086$ & $+0.0536$ & $+0.1000$ \\
$B$ & $+0.2000$ & $+0.0670$ & $-0.0403$ & $-0.2487$ \\
$\varepsilon$ & $+0.0028$ & $+0.0416$ & $-0.0133$ & $+0.1487$ \\
\end{tabular}
\end{center}

Every column sums to $0$ — mass pushed off the truth's cells lands
somewhere — and on this example even the printed four-decimal digits
sum to $0.0000$ exactly. Two cells are worth reading. At $t = 1$,
$B$'s entry is \emph{exactly} $+0.2000 = y$: the state is
unreachable, its occupancy is $0$, and the gradient is the output
itself — pure push-down. And scoring the \emph{wrong} transcript BA
against the same matrix flips $t{=}1$, $A$ to exactly $+0.7000$ for
the same reason: an alignment the transcript cannot use gets no
protection.

The degeneration closes the derivative chapter's promise: let one
path carry everything and $\gamma$'s columns snap one-hot —
$y - \gamma$ becomes $p - \text{one-hot}$, the
$\anchor{009.G.nllgradient}$ of the toolkit's closer, received as
the special case. Bridle's ``output minus a one-from-N
target''~\cite{probability-john-s-bridle-training-stochastic-model-recognition-nips-198},
with the one-from-N gone soft. And one implementation trap, because
the identity is what CTC backward passes hard-code: the clean form
lives at the \emph{logits}. Hand the loss anything but a true
log-softmax and the forward loss stays right while the gradient goes
silently
wrong~\cite{deep_learning-pytorch-issue-122243-non-normalized-ctc-inputs}.
(The $y$-space form $-\gamma/y$ is all-negative — it only pushes up —
which is why the $u$-form is the teachable one.)

\section{The error signal learns}

Watch the identity during training. Graves's original figure traced
the error signal through three
stages~\cite{deep_learning-graves-et-al-2006-connectionist-temporal-classification};
here the arc is rebuilt countable, by plain gradient descent — the
previous chapter's walker, learning rate $1.0$ — on free per-frame
logits. (One sign note where plots enter: error panels traditionally
plot the \emph{push} $\gamma - y$, the negative of the gradient —
above the axis means ``raise this output''. Columns balance to zero
in either sign.)

\emph{Diffuse.} Start uninformative: uniform outputs. Then $\gamma$
is exactly path counts over $15$ — the input never consulted — and
the gradient is pure fractions: at $t = 1$,
$y - \gamma = \anchor{010.L.gradt1}$ across $(A, B, \varepsilon)$.
The error is determined by the target sequence alone, spread wide
and time-symmetric.

\emph{Localised.} Descend from the worked matrix: the loss walks
$\anchor{010.K.NLL} \to \anchor{010.M.loss10}$ after ten steps
$\to \anchor{010.M.loss50}$ after fifty. The error concentrates
where the outputs already lean — the pushes sharpen around the
emerging alignment $A, A, \varepsilon, B$.

\emph{Gone.} By five thousand steps the loss reads
$\anchor{010.M.loss5000}$ and the bars sink into the axis. The
teaching gem hides at frame $3$, which converges not to a spike but
to the \emph{mixed} output $\anchor{010.M.frame3}$ with gradient
$\to 0$: once frames $1, 2, 4$ say $A, A, B$, all three of frame
$3$'s choices collapse to AB, so the loss goes indifferent — $y$
matches $\gamma$ out of indifference, not certainty. A vanished
gradient does not mean one-hot outputs; it means the outputs agree
with the occupancy, which is a weaker and more interesting thing.

\section{Why the spikes appear}

Trained CTC models emit sharp single-frame spikes separated by long
runs of blank, and the alignment chapter warned: never read the
spikes as segment boundaries. Here, finally, is the mechanism — and
it is topology plus weight sharing, not acoustics.

First, topology. Count label occurrences over all AB-paths, with the
input nowhere in the computation. At $T = 4$ the counts across
$(A, B, \varepsilon)$ are $\anchor{010.N.countsT4}$ — blank is
\emph{not} dominant; the boundary case is honest. At $T = 5$ (now
$\anchor{010.N.t5paths}$ paths — the count the alignment chapter's
problem set met) they are $\anchor{010.N.countsT5}$, and blank is
ahead for good: its share grows with $T$, toward $2/3$ in the
one-letter limit under uniform outputs. Under uniform initialization
the occupancy \emph{is} these counts — so the earliest training
signal already tilts blankward for every $T$ past the boundary.

Second, weight sharing. Free per-frame outputs never spike — the
worked example converged to an honest alignment above. But force one
shared softmax to serve every frame (a bias-only model) and descent
tells a different story: at $T = 4$ it settles at the honest
$\anchor{010.N.biasT4}$; at $T = 12$ it converges to
$\anchor{010.N.biasT12}$ — argmax blank at every frame, an
\emph{empty} decode, $100\%$ error — and this is a \emph{local}
optimum: the global optimum has zero error. That trap is proved, for
a feed-forward network from uniform initialization, by Zeyer,
Schl\"uter and Ney (arXiv
preprint)~\cite{deep_learning-zeyer-et-al-2021-why-does-ctc-result-in-peaky-behavior},
who also name the repair: a label prior folded into the loss —
peakiness is a steerable training artifact, never a timestamp.

The closing map: the identity is a family. One-hot cross-entropy
($\gamma$ one-hot — the degenerate case), distillation (a teacher's
soft targets), CTC ($\gamma$ the occupancy) — all one gradient
shape, output minus target, with only the target changing. And
$\gamma$ itself is the soft alignment real systems compute — the
E-step object of the HMM tradition, met here as the thing the
gradient reads off the trellis.

\paragraph{Where this goes.} Training is solved; reading the trained
model's output stream is not. Greedy decoding has been suspect since
the alignment chapter, spikes are not segmentations, and the honest
decoder — beam search over \emph{collapsed prefixes}, each prefix
carrying two probabilities — is the road's last construction: the
next chapter closes the loop from training to deployment.
\end{narrative}

\section*{Practice problems}

\begin{problem}
At unit weights, run the backward recurrence on the AB trellis
($T = 4$, states $\varepsilon, A, \varepsilon, B, \varepsilon$):
$\beta_4 = 1$ at the two accepting states, $0$ elsewhere. Compute
the $t = 1$ column of suffix counts, and check its two initial
states against the flagship path count.
\begin{hint}
Work right to left; a state's count is the sum of its successors'
counts (stay, advance, and the skip only over a blank between
different letters).
\end{hint}
\begin{solution}
From $\beta_4 = (0, 0, 0, 1, 1)$: the $t = 3$ column is
$(0, 1, 1, 2, 1)$ — e.g.\ state $B$ has stay $1$ + advance $1$; the
$A$ state's skip to $B$ is legal since $A \neq B$. Then $t = 2$
gives $(1, 4, 3, 3, 1)$ — $\beta_2(A) = \text{stay } 1 +
\text{advance } 1 + \text{skip } 2 = 4$ — and $t = 1$ gives
$(5, 10, 6, 4, 1)$. The two initial states (leading blank and first
letter): $5 + 10 = 15$, the same count the forward sweep found from
the other end.
\end{solution}
\end{problem}

\begin{problem}
Using the forward counts $(1,1,0,0,0)$, $(1,2,1,1,0)$,
$(1,3,3,4,1)$, $(1,4,6,10,5)$ and your backward counts, form the
per-cell products $\alpha \beta$ at every cell and sum each column.
What do the four sums say, and why must they agree?
\begin{hint}
Backward columns: $t{=}1$ $(5,10,6,4,1)$, $t{=}2$ $(1,4,3,3,1)$,
$t{=}3$ $(0,1,1,2,1)$, $t{=}4$ $(0,0,0,1,1)$.
\end{hint}
\begin{solution}
The product columns are $(5, 10, 0, 0, 0)$, $(1, 8, 3, 3, 0)$,
$(0, 3, 3, 8, 1)$ and $(0, 0, 0, 10, 5)$ — each summing to $15$.
$\alpha_t(s) \beta_t(s)$ counts the paths \emph{through} cell
$(t, s)$ — ways in times ways out — and since every one of the $15$
paths occupies exactly one state at every frame, summing over a
column counts every path exactly once, whichever column you pick.
(The $8$ at $(t{=}2, A)$ is the waist: $2$ prefixes $\times$ $4$
suffixes.) On the real matrix the same sweep returns
$P(\mathrm{AB} \mid X) = 0.4877$ four times — the constant column,
with probabilities riding the counts.
\end{solution}
\end{problem}

\begin{problem}
Under uniform outputs ($y_t(k) = 1/3$ for all $t, k$): compute
$P(\mathrm{AB} \mid X)$, the occupancy column at $t = 1$, and the
gradient $y - \gamma$ there. What is remarkable about what the
computation never used?
\begin{hint}
Every path has the same weight, so $\gamma$ is a ratio of counts.
\end{hint}
\begin{solution}
Each of the $15$ paths weighs $(1/3)^4 = 1/81$, so
$P = 15/81 = 5/27$. With all weights equal, occupancy is path
counts over $15$: at $t = 1$, ten paths start in $A$ and five in the
leading blank, so $\gamma_1 = (2/3, 0, 1/3)$ across
$(A, B, \varepsilon)$. The gradient is
$y - \gamma = (1/3 - 2/3,\ 1/3 - 0,\ 1/3 - 1/3) =
(-1/3,\ +1/3,\ 0)$. The input never entered: at an uninformative
start, the error signal is determined by the target sequence's
topology alone — which is exactly the diffuse first stage of
training, and the seed of blank's topological head start at larger
$T$.
\end{solution}
\end{problem}

\begin{problem}
On the worked matrix, frame $1$ has $y_1 = (0.7, 0.2, 0.1)$ and
occupancy $\gamma_1 = (0.9028, 0, 0.0972)$. Write the gradient
column, verify it sums to zero, and explain why $B$'s entry equals
its softmax output exactly — no rounding involved.
\begin{hint}
Which trellis states can a path for AB occupy at frame $1$?
\end{hint}
\begin{solution}
$y - \gamma = (0.7 - 0.9028,\ 0.2 - 0,\ 0.1 - 0.0972) =
\anchor{010.K.gradt1}$, summing to $0$: both $y$ and
$\gamma$ sum to $1$ over the classes, so their difference sums to
zero — mass pushed off the truth's cells must land on the others. At
frame $1$ a path for AB can occupy only the leading blank or the
first $A$; the $B$ state is unreachable, so its occupancy is exactly
$0$ and its gradient is exactly $y = +0.2000$ — the loss pushes a
provably useless activation straight down. (The same mechanism, run
against the wrong transcript BA, flips frame $1$'s $A$ to exactly
$+0.7000$.)
\end{solution}
\end{problem}

\begin{problem}
Derive the identity. Starting from
$\partial L / \partial \ln y_t(k) = -\gamma_t(k)$ and the
log-softmax derivative
$\partial \ln y_j / \partial u_k = \delta_{jk} - y_k$, show
$\partial L / \partial u_t(k) = y_t(k) - \gamma_t(k)$, naming the
step where the constant column does the work.
\begin{hint}
Chain over $j$, then split the sum into its $\delta$ term and its
$y_k$ term.
\end{hint}
\begin{solution}
Chaining,
\[
  \frac{\partial L}{\partial u_t(k)}
  = \sum_j \frac{\partial L}{\partial \ln y_t(j)}
    \frac{\partial \ln y_t(j)}{\partial u_t(k)}
  = -\sum_j \gamma_t(j) \bigl(\delta_{jk} - y_t(k)\bigr) .
\]
The $\delta$ term picks out $-\gamma_t(k)$; the second term is
$+\,y_t(k) \sum_j \gamma_t(j)$, and the sum is $1$ — the column
checksum, the same constant column that verified the bookkeeping,
now killing the Jacobian's second term. What remains is
$y_t(k) - \gamma_t(k)$: softmax output minus occupancy. With a
single surviving path, $\gamma$ is one-hot and the identity is the
derivative chapter's $p - \text{one-hot}$ — the general case
received the special one.
\end{solution}
\end{problem}

\begin{problem}[stretch]
Count label cells over all AB-paths — with no input anywhere — at
$T = 4$ and $T = 5$: how many of the path-cells carry $A$, $B$, and
blank? At which $T$ does blank take the lead, and what does that say
about the first gradient step from an uninformative start?
\begin{hint}
$15$ paths of $4$ cells, then $35$ paths of $5$ cells; enumerate or
use the unit-weight $\gamma$ columns, which are exactly these counts
over the path total.
\end{hint}
\begin{solution}
At $T = 4$: $60$ path-cells split $(A, B, \varepsilon) =
(21, 21, 18)$ — the labels lead, and blank is \emph{not} dominant;
any unconditioned ``blank always dominates'' claim dies on this
case. At $T = 5$: $175$ path-cells split $(56, 56, 63)$ — blank
ahead, and it stays ahead for every larger $T$, with its share
growing toward $2/3$ in the one-letter limit under uniform outputs.
Since uniform-init occupancy equals these count shares, the very
first descent step already pushes blankward once $T \geq 5$ — the
topological half of peakiness, before weight sharing adds the trap.
\end{solution}
\end{problem}
